% Sections 2.1 to 2.7, from lectures/L02/appendix/a_larger_networks.md: one section per part,
% A.1 to A.7.

\section{Karnaugh Maps: From Truth Table to Minimized Gate Network}
\label{c2:a1}
Reading a sum-of-products equation straight off a truth table, the way \secref{c1:a2} did, always
works but rarely gives the smallest network: every \code{1} row becomes its own AND-term, even when
several rows share most of their inputs and could be combined.

A \textbf{Karnaugh map} (K-map) makes those shared patterns visible by hand, for functions of 2 to 4
variables, with no formal Boolean-algebra manipulation. One worked example follows; further practice
is in \secref{c2:sec:exercises}, and classical minimization theory (don't-cares, Quine--McCluskey)
is out of scope.

The idea is to lay the truth table out as a grid rather than a list, ordering the header labels so
every cell differs from each horizontal and vertical neighbour in exactly one input value, including
around the left/right and top/bottom edges. That ordering is \textbf{Gray code}: for two bits,
\code{00, 01, 11, 10}, where only one bit changes at each step, including from the last entry back
to the first.

Then:
\begin{enumerate}
  \item Write a \code{1} in every cell where the output is \code{1}, and leave the rest blank.
  \item Draw rectangular groups around the \code{1}s:
    \begin{itemize}
      \item Every group must be a rectangle whose width and height are each a power of 2 (1, 2 or
        4), so its size is a power of 2 as well.
        \begin{itemize}
          \item a group of 4 is therefore either a $1\times4$ line or a $2\times2$ block, and both
            are allowed.
          \item cells that touch only at a corner are not a rectangle, so they cannot be grouped.
        \end{itemize}
      \item Groups may overlap, and may wrap around an edge of the grid.
    \end{itemize}
  \item Make every group as large as possible, and cover every \code{1} at least once. A group of 1
    keeps all its inputs in the term, a group of 2 drops one, a group of 4 drops two, and so on.
    \begin{itemize}
      \item Then use as few groups as will do it: drop any group whose cells are every one of them
        already covered by another. Nothing above forbids a redundant group, and a redundant group
        is a redundant AND-term in the answer, correct but not minimal.
    \end{itemize}
  \item For each group, keep only the input(s) that hold the same value in every cell of that
    group, dropping the rest, and write that as one AND-term:
    \begin{itemize}
      \item write a kept input directly if that constant value is \code{1}, and primed if it is
        \code{0}.
      \item so a group covering only cells where \code{B = 0} contributes the term \code{B'}.
    \end{itemize}
  \item \code{X} is the OR of all the group terms.
\end{enumerate}

\subsection{Worked example}
\label{c2:kmapexample}
Three inputs \code{A}, \code{B}, \code{C}, one output \code{X}:

\begin{narrowtable}{@{}cccc@{}}
  \thead{A} & \thead{B} & \thead{C} & \thead{X} \\
  \midrule
  0 & 0 & 0 & 0 \\
  0 & 0 & 1 & 1 \\
  0 & 1 & 0 & 0 \\
  0 & 1 & 1 & 1 \\
  1 & 0 & 0 & 0 \\
  1 & 0 & 1 & 1 \\
  1 & 1 & 0 & 1 \\
  1 & 1 & 1 & 1 \\
\end{narrowtable}

\noindent Lay it out with \code{AB} in Gray code down the rows and \code{C} across the columns,
filling in a \code{1} wherever \code{X = 1} (Figure~\ref{c2:fig:karnaugh_filled}).

\bookfigure{karnaugh_filled}{The Karnaugh map for \code{X}, filled in from the truth table:
\code{AB} in Gray code down the rows, \code{C} across the columns.}{c2:fig:karnaugh_filled}

\noindent The entire \code{C = 1} column is four \code{1}s, a group of 4. Every cell shares
\code{C = 1} and nothing else is constant (\code{A} and \code{B} both vary), so it reduces to the
single-variable term \code{C}:

\bookfigure{karnaugh_group_c}{The same map with the four-cell \code{C = 1} group marked, which
reduces to the term \code{C}.}{c2:fig:karnaugh_group_c}

\noindent One \code{1} is still uncovered, on the \code{AB = 11} row. Group it with its neighbour to
the right, which the \code{C} group already covers. Every cell in the pair shares \code{A = 1} and
\code{B = 1}, and \code{C} varies, so it reduces to \code{AB}:

\bookfigure{karnaugh_groups}{The same map with the two-cell \code{AB = 11} group added, overlapping
the \code{C} group in one cell.}{c2:fig:karnaugh_groups}

\noindent Every \code{1} is now covered; the \code{AB = 11, C = 1} cell is in both groups, and
overlap is fine. Summing gives:

\begin{console}
X = AB + C
\end{console}

\noindent Reading the truth table directly, as Chapter~\ref{c1:ch} did, would have given five
AND-terms, one per \code{1} row. The K-map found an equivalent two-term equation, which is the entire
value of doing this by hand.

\subsection{Realizing the network}
\label{c2:kmapnetwork}
\code{X = AB + C} is two gates: an \code{AND} computing \code{AB}, feeding an \code{OR} whose other
input is \code{C}.

\bookimage{0.55}{lectures/L02/appendix/images/net.png}{The gate network for the minimized
expression \code{X = AB + C}: an \code{AND} gate feeding an \code{OR}, built in
CircuitVerse.}{c2:fig:net}

\noindent The course builds this network by hand in CircuitVerse, live; build it yourself exactly as
you built \code{or_gate}'s in Chapter~\ref{c1:ch}. \Secref{c2:a2} translates it into VHDL.

\section{Translating the Karnaugh-Map Network into VHDL}
\label{c2:a2}
\code{X = AB + C} translates straight into VHDL. Three inputs, one output:

\begin{vhdlcode}
library ieee;
use ieee.std_logic_1164.all;

entity gate_network is
    port(a, b, c: in std_logic;
         x      : out std_logic);
end entity;
\end{vhdlcode}

\noindent The architecture uses an internal signal \code{ab} to hold the \code{AND} gate's output,
mirroring the two-gate network one-for-one:

\needspace{12\baselineskip}
\begin{vhdlcode}
architecture behaviour of gate_network is
signal ab: std_logic;
begin
    -- AND gate: AB.
    ab <= a and b;

    -- OR gate: AB + C.
    x <= ab or c;
end architecture;
\end{vhdlcode}

\noindent Without the intermediate signal this collapses to \code{x <= (a and b) or c;}. Both are
equally valid and synthesize to the same hardware; use whichever is more readable for the network at
hand.

The buildable version, \file{lectures/L02/gate_network/gate_network.vhd}, uses the collapsed form,
since at two gates there is nothing to gain from naming the intermediate. It is printed in full
below; set it beside the version above and satisfy yourself they describe the same circuit. That is
the claim this section is making, and it costs nothing to check.

\vhdlfile{lectures/L02/gate_network/gate_network.vhd}

To run it on hardware, assign the ports to switches and an LED with the DE0-CV's pin planner
(Appendix~\ref{app:quartus}). Nothing here depends on specific pin numbers.

\section{\code{std_logic_vector}: More Than One Wire}
\label{c2:a3}
Everything so far has been \code{std_logic}, one wire carrying one bit. Bundling several wires under
one name is what \code{std_logic_vector} is for.

\begin{vhdlcode}
signal count: std_logic_vector(3 downto 0);
\end{vhdlcode}

\noindent That declares four wires: \code{count(3)} down to \code{count(0)}. The range is
\code{3 downto 0} rather than \code{0 to 3} so the leftmost bit is the most significant, the way you
would write the number on paper. This book uses \code{downto} everywhere, and so does almost all the
VHDL you will meet.

A vector is a port type like any other:

\begin{vhdlcode}
entity display is
    port(number: in  std_logic_vector(3 downto 0);
         hex   : out std_logic_vector(6 downto 0));
end entity;
\end{vhdlcode}

\subsection{Values}
\label{c2:values}
A single bit takes single quotes, a vector double quotes:

\begin{vhdlcode}
x     <= '1';                -- std_logic: one bit
count <= "1010";             -- std_logic_vector(3 downto 0): four bits
count <= (others => '0');    -- every bit '0', whatever the width
\end{vhdlcode}

\noindent The last is an \emph{aggregate}: \code{(others => '0')} means ``all the remaining bits are
\code{'0'}'', and since none were named individually, that is all of them. It keeps working when the
width changes, and it appears throughout the rest of the book.

\subsection{Indexing and slicing}
\label{c2:slicing}
One bit out of a vector is a \code{std_logic}; a contiguous run is a \emph{slice}, itself a vector:

\begin{vhdlcode}
carry     <= count(3);
high_half <= input(7 downto 4);
\end{vhdlcode}

\noindent Slicing is how one wide port feeds several narrower ones, which is what \secref{c2:a6}'s
two-digit display does with its eight switch inputs.

One detail is not what a software engineer expects. \code{input(7 downto 4)} has the index range
\code{7 downto 4}, not \code{3 downto 0}. Connected to a port declared
\code{std_logic_vector(3 downto 0)}, the two are matched \textbf{by position, left to right}, not by
index number: \code{input(7)} drives bit 3, \code{input(6)} drives bit 2, and so on. The lengths
have to match. The numbering does not.

\subsection{A vector is a bundle of bits, not a number}
\label{c2:bundle}
\code{count <= "1010"} does not make \code{count} equal ten. \code{std_logic_vector} says nothing
about what the bits \emph{mean}, so it offers only the operations that make sense for a bundle of
wires: the logical operators applied bit by bit (\code{a and b}, \code{not a}, \code{a xor b}),
comparison (\code{count = "1010"}), and concatenation (\code{"10" & "11"} is \code{"1011"}).

Arithmetic is not on that list. \code{count + 1} will not compile, because nothing has said whether
those four wires are an unsigned number, a signed one, or four unrelated signals.

This book does not work around that with a numeric vector type. Anything counted is a plain integer
with a range:

\begin{vhdlcode}
signal count: natural range 0 to 9;
\end{vhdlcode}

\noindent which you can add to and compare directly, and which synthesis turns into exactly the
flip-flops the range needs. Vectors stay what they are: bundles of wires, for things that really are
bundles of wires, like eight switches or seven display segments.

The two worlds meet only at a boundary, and the conversions appear where they are needed:

\begin{vhdlcode}
count_v <= std_logic_vector(to_unsigned(count, 4));   -- integer to vector
count   <= to_integer(unsigned(count_v));             -- vector to integer
\end{vhdlcode}

\noindent Both come from \code{ieee.numeric_std}, which is why some files add a third \code{use}
clause.

\section{Multiplexers}
\label{c2:a4}
A \textbf{multiplexer} (``mux'') routes exactly one of several data inputs to a single output, chosen
by a separate set of \textbf{selector} inputs. An \code{N}-to-1 mux has \code{N} data inputs and
\code{ceil(log2(N))} selector bits.

The simplest case, a 2-to-1 mux with inputs \code{A}, \code{B}, selector \code{S} and output
\code{X}:

\begin{narrowtable}{@{}cc@{}}
  \thead{S} & \thead{X} \\
  \midrule
  0 & B \\
  1 & A \\
\end{narrowtable}

\noindent As a Boolean equation, \code{X = S'B + SA}: for every input combination exactly one
AND-term is active, and that term's data input passes through.

Larger muxes are the same idea with more selector bits, though \emph{which} data input a given
selector value names is a convention that changes from part to part, and has to be read off the
specification rather than assumed. The 8-to-1 below is the opposite way round from the 2-to-1
above: there \code{0} names the second input, here \code{000} names the first. An 8-to-1 mux with
data inputs \code{A}--\code{H} and selector \code{S[2:0]}:

\begin{narrowtable}{@{}cc@{}}
  \thead{S[2:0]} & \thead{X} \\
  \midrule
  000 & A \\
  001 & B \\
  010 & C \\
  011 & D \\
  100 & E \\
  101 & F \\
  110 & G \\
  111 & H \\
\end{narrowtable}

\noindent Read into sum-of-products one term per row, the \code{000} row contributes
\code{A*S2'S1'S0'} and the rest follow the same shape, each gating its data input with the selector
combination that selects it:

\begin{console}
X = AS2'S1'S0' + BS2'S1'S0 + ... + HS2S1S0
\end{console}

\noindent Worth writing out once by hand, not worth memorizing. The pattern is the point, and it is
the same at any size.

As gates: one 4-input \code{AND} per data input (the data bit plus all three selector bits),
inverting selector bits as needed, with all eight \code{AND} outputs summed by a tree of \code{OR}
gates (Figure~\ref{c2:fig:mux8}).

\bookimage[htbp]{0.85}{lectures/L02/appendix/images/mux_8to1.png}{The 8-to-1 multiplexer as gates: eight
4-input \code{AND} gates, each gating one data input with the selector combination that selects it,
feeding a tree of \code{OR} gates.}{c2:fig:mux8}

\noindent This is a real building block. The ATmega328P (the chip on an Arduino Uno/Nano) uses one
internally so its analog-capable pins, plus an internal bandgap reference and a temperature sensor,
can share a single analog-to-digital converter, selected by the four \code{MUX} bits of its
\code{ADMUX} register.

\section{Multiplexers in VHDL: \code{process} and \code{case}}
\label{c2:a5}
Every statement so far has been \textbf{concurrent}: a standing description of a wire, all active at
once. A multiplexer is the first circuit that is awkward to write that way, and needs two new
constructs:
\begin{itemize}
  \item a \textbf{\code{process}}, a block whose contents run sequentially, top to bottom, like
    ordinary code. It re-runs whenever any signal in its \textbf{sensitivity list} changes, and
    from the architecture's point of view it is still just one more concurrent statement. Every
    process in this book has this shape; Chapter~\ref{c3:ch} reuses it for clocked logic, with the
    clock in the sensitivity list instead of the inputs being decoded.
  \item a \textbf{\code{case} statement}, VHDL's \code{switch}. Like \code{if} and \code{for} it is
    a \emph{sequential} statement, legal only inside a process (or a subprogram, which this book
    does not use), and every value the selector can take must be covered, which
    \code{when others} guarantees.
\end{itemize}
``Pick one of several inputs based on a selector'' \emph{is} a \code{case} statement, almost word for
word. The 2-to-1 mux from \secref{c2:a4}:

\needspace{11\baselineskip}
\begin{vhdlcode}
MUX_PROCESS: process(a, b, sel) is
begin
    case (sel) is
        when '0'    => x <= b;
        when '1'    => x <= a;
        when others => x <= '0';
    end case;
end process;
\end{vhdlcode}

\noindent Two points from that:
\begin{itemize}
  \item Single-bit values take single quotes, \code{'0'}/\code{'1'}, as opposed to the double quotes
    used for vectors like \code{"000"}.
  \item \code{when others} is required rather than defensive. VHDL insists a \code{case} cover every
    value of its selector's type, and \code{std_logic} has nine (\secref{c1:a4}), so \code{'0'} and
    \code{'1'} alone never make one exhaustive. This is a language rule, not a claim about the wire:
    a synthesized net carries a \code{0} or a \code{1} and nothing else, and the other seven values
    are the simulator's, where a signal sitting at \code{'U'} or \code{'X'} is a bug rather than a
    level.
\end{itemize}

\subsection{Worked example: the 8-to-1 mux}
\label{c2:mux8}
Same shape with a 3-bit selector and eight branches, one per row of \secref{c2:a4}'s truth table.
The data inputs \code{A}--\code{H} become the vector \code{inputs(7 downto 0)}, with \code{A} as
\code{inputs(7)} down to \code{H} as \code{inputs(0)}, so selector \code{"000"} routes
\code{inputs(7)} and \code{"111"} routes \code{inputs(0)}.

The buildable module is \file{lectures/L02/mux_8to1/mux_8to1.vhd}:

\vhdlfile{lectures/L02/mux_8to1/mux_8to1.vhd}

Compare it against \secref{c2:a4}'s eight-term equation for the same circuit. The equation describes
the gates; the \code{case} describes the \emph{intent} and lets the synthesis tool produce the gates.
Both are correct and both synthesize to the same hardware. That gap, between describing structure
and describing behaviour, is most of what a hardware description language buys you.

\section{Building a Design from Submodules}
\label{c2:a6}
Everything so far has been one entity with one architecture. That stops being enough the moment a
design needs the same block twice.

The DE0-CV has six 7-segment displays, and you want to drive two, each showing one hexadecimal
digit. The logic turning a 4-bit number into seven segment signals is the same for both. You could
write it twice inside one architecture and nothing would stop you, but the second copy is a
liability: every fix has to be made in both, and the day they drift apart is the day you lose an
evening working out which one is wrong.

Write the block once, as its own entity, and \emph{instantiate} it twice.

\subsection{The submodule}
\label{c2:submodule}
An ordinary module. Nothing about it says ``I am a submodule'':

\begin{vhdlcode}
entity display is
    port(number: in  std_logic_vector(3 downto 0);
         hex   : out std_logic_vector(6 downto 0));
end entity;
\end{vhdlcode}

\noindent Whether it ends up as a whole design or one piece of a larger one is decided by whoever
instantiates it, not by the module.

\subsection{Instantiating it}
\label{c2:instantiate}
With \code{entity work.<name>}, inside another architecture:

\begin{vhdlcode}
entity hex_display is
    port(input     : in  std_logic_vector(7 downto 0);
         hex1, hex0: out std_logic_vector(6 downto 0));
end entity;

architecture structure of hex_display is
begin
    display1: entity work.display
        port map(input(7 downto 4), hex1);

    display0: entity work.display
        port map(input(3 downto 0), hex0);
end architecture;
\end{vhdlcode}

\begin{itemize}
  \item \code{display1} and \code{display0} are \textbf{instance labels}, naming the two copies.
    They must be unique within the architecture, and they are not the module's name: both instances
    are \code{display}.
  \item \code{entity work.display} selects the module. \code{work} is the library your analyzed
    files land in by default, in GHDL and Quartus alike.
  \item \code{port map(...)} connects the instance's ports to signals in the enclosing architecture,
    \textbf{in the order that entity declares them}. \code{display} declares \code{number} then
    \code{hex}, so the first item drives \code{number} and the second is driven by \code{hex}.
  \item \code{input(7 downto 4)} is a \textbf{slice}, which the instance sees as a
    \code{std_logic_vector(3 downto 0)}.
\end{itemize}
Note what \code{hex_display}'s architecture does \emph{not} contain: any logic. It is wiring and
nothing else. Every gate in the finished design comes from the two \code{display} instances.

\subsection{Positional and named association}
\label{c2:association}
The \code{port map} above is \textbf{positional}, matching ports by order. VHDL also allows
\textbf{named} association, which you will meet in other people's code:

\begin{vhdlcode}
display1: entity work.display
    port map(number => input(7 downto 4), hex => hex1);
\end{vhdlcode}

\noindent The two are equivalent. This book uses positional everywhere, which means an entity's port
declaration order is something you have to get right. That is not a formality: it is exactly what
the testbenches you are given rely on, which is why every exercise's self-check line spells out its
ports \emph{in order}.

\subsection{What this costs in hardware}
\label{c2:cost}
Nothing surprising. Instantiating a module twice puts two copies of its gates on the FPGA, exactly
as writing the logic twice would. The saving is in the source, not the silicon: one description, one
place to fix, and a design you can read.

\section{Checking Your Work}
\label{c2:a7}
Every VHDL module you write for an exercise comes with a \textbf{self-checking testbench}: a small
VHDL file that drives your inputs, checks your outputs, and tells you which case failed.
Chapter~\ref{c1:ch}'s Exercise~\ref{c1:ex:2} already shipped one, and \secref{c1:a5} gave you the
three commands that run it.

You do not write testbenches in this book. Running them is enough to verify every exercise on your
own laptop, with no FPGA board.

\Secref{c2:app:c} is the full guide: what a testbench is and why its entity has no ports, the
analyze $\to$ elaborate $\to$ run sequence, how to read a failed assertion, and how your module's
\textbf{port order} has to match what the testbench expects.

Run one now, against this chapter's worked example:

\begin{shell}
cd lectures/L02/gate_network
ghdl -a --std=93 gate_network.vhd gate_network_tb.vhd
ghdl -e --std=93 gate_network_tb
ghdl -r --std=93 gate_network_tb --assert-level=error --stop-time=10ms
\end{shell}
